% ====================================================================
\chapter{Final Master Review}\label{ch:review}
% ====================================================================
\markboth{Final Master Review}{}

This final part draws the whole book together. It is designed for revision: read the module summary and the concept map
to see the structure, use the formula sheet and the glossary as references, work through the checklist honestly, and then
attempt the final practice paper under timed conditions (about $1.2$ minutes per mark, without a calculator) before reading
the solutions.

\section*{Module summary}
\addcontentsline{toc}{section}{Module summary}

\paragraph{The language (\cref{ch:logic,ch:sets})} Mathematics is written in propositions --- statements that are true or
false --- linked by implication, equivalence and negation, and quantified with ``for all'' ($\forall$) and ``there exists''
($\exists$). Axioms are accepted without proof; theorems are proved from them. Sets collect objects; they are compared by
$\subseteq$ and $=$ (two sets are equal exactly when each is a subset of the other) and combined by $\cap$, $\cup$ and
$\setminus$. Set-builder notation $\{x\in A\mid P\}$ describes a set by a property, and solving an inequality amounts to
describing such a set simply.

\paragraph{The numbers (\cref{ch:numbers})} $\N\subseteq\Z\subseteq\Q\subseteq\R$. The real numbers satisfy thirteen
properties: trichotomy and transitivity of the order (O1--O2), five properties of addition (A1--A5) and six of
multiplication (M1--M6). The rationals satisfy them too; the real numbers are distinguished by the (non-examinable) Least
Upper Bound property, which says the number line has no gaps. Intervals name the sets of numbers between two endpoints.

\paragraph{Calculation (\cref{ch:ineq,ch:abs,ch:sum})} Inequalities are solved by chains of equivalences: adding is always
safe, multiplying by a positive number keeps the direction, by a negative number reverses it (M6). Products are handled by
sign analysis; quadratics are factorised through their roots using the discriminant; fractions are handled by bringing
everything to one side. The absolute value measures size and $|x-y|$ measures distance; the triangle inequality says a
detour is never shorter than the direct route. The summation operator $\sum$ writes long sums compactly and obeys P1--P3,
with the conventions that a one-term sum is that term and an empty sum is $0$.

\paragraph{Proof and structure (\cref{ch:ind,ch:binom,ch:func})} Induction proves a proposition for every natural number
by proving a base and a step. It establishes summation formulas, the binomial formula
$(x+y)^n=\sum_jC_n^jx^{n-j}y^j$, and --- in its strong form --- the factorisation $p(x)=(x-\bar x)q(x)$ of a polynomial
through any root. Functions associate each element of a domain with a unique element of a codomain; two functions are equal
only if domains, codomains and values agree.

\paragraph{Limits (\cref{ch:limits})} A sequence is a function on $\N$. A number $x$ is its limit if for every
$\varepsilon>0$ there is $N$ such that $|a_n-x|<\varepsilon$ for all $n>N$. Convergence is proved by producing $N$ for an
arbitrary $\varepsilon$; non-convergence by contradiction, using the triangle inequality. Limits are unique and preserve
non-strict inequalities. This definition is the gateway to the rest of the module.

\section*{How the ideas connect}
\addcontentsline{toc}{section}{How the ideas connect}

\begin{figure}[ht]
\centering
\begin{tikzpicture}[node distance=8mm and 9mm, every node/.style={mapnode=34mm}]
  \node (prop) {Propositions,\\ $\Rightarrow$, $\Leftrightarrow$, not};
  \node[right=of prop] (quant) {Quantifiers $\forall$, $\exists$};
  \node[right=of quant] (proof) {Proof: direct, cases,\\ contradiction};
  \node[below=of prop] (sets) {Sets: $\subseteq$, $\cap$, $\cup$, $\setminus$};
  \node[below=of quant] (axioms) {Axioms of $\R$:\\ O1--2, A1--5, M1--6};
  \node[below=of proof] (ind) {Induction\\ (successor property)};
  \node[below=of sets] (int) {Intervals and\\ solution sets};
  \node[below=of axioms] (ineq) {Inequalities: M6,\\ sign analysis, $D$};
  \node[below=of ind] (sum) {Summation $\sum$,\\ P1--P3};
  \node[below=of int] (func) {Functions, polynomials,\\ factor theorem};
  \node[below=of ineq] (abs) {Absolute value,\\ triangle inequality};
  \node[below=of sum] (binom) {Binomial formula};
  \node[below=of abs, mapnode=70mm, fill=corebg, draw=corecol] (lim) {Sequences and the limit:\\
     $\forall\varepsilon>0\ \exists N\ \forall n>N:\ |a_n-x|<\varepsilon$};
  \draw[maparrow] (prop) -- (quant); \draw[maparrow] (quant) -- (proof);
  \draw[maparrow] (prop) -- (sets); \draw[maparrow] (quant) -- (axioms); \draw[maparrow] (proof) -- (ind);
  \draw[maparrow] (sets) -- (int); \draw[maparrow] (axioms) -- (ineq); \draw[maparrow] (ind) -- (sum);
  \draw[maparrow] (int) -- (ineq);
  \draw[maparrow] (int) -- (func); \draw[maparrow] (ineq) -- (abs); \draw[maparrow] (sum) -- (binom);
  \draw[maparrow] (func) |- (lim); \draw[maparrow] (abs) -- (lim); \draw[maparrow] (binom) |- (lim);
\end{tikzpicture}
\caption{Concept map of the book. Arrows point from ideas to the ideas that use them. Every route ends at the definition of
the limit, which needs quantifiers, inequalities, distance and induction together. (Induction is also used to prove the
factor theorem for polynomials, \cref{thm:factor}.)}
\label{fig:concept-map}
\end{figure}

\paragraph{Five threads that run through everything}
\begin{enumerate}[label=\arabic*.]
  \item \emph{Equivalence preserves solution sets.} Whether solving an inequality (\cref{ch:ineq}) or proving set equality
  (\cref{ch:sets}), the goal is a two-way statement.
  \item \emph{The sign of a multiplier.} M6 drives linear inequalities, sign analysis, parameter cases, quadratic and
  fractional inequalities, and every manipulation of inequalities inside limit proofs.
  \item \emph{Arbitrary versus particular.} $\forall$ statements are proved for an arbitrary object, $\exists$ statements by
  an example; this structures subset proofs, induction steps and limit proofs alike.
  \item \emph{Distance.} $|x-y|$ and the triangle inequality turn ``close to'' into algebra.
  \item \emph{Building on the previous case.} Induction, the summation property P3 (peeling off the last term) and Pascal's
  rule all express a structure in terms of a smaller one.
\end{enumerate}

\clearpage
\section*{Master formula sheet}
\addcontentsline{toc}{section}{Master formula sheet}

{\small
\renewcommand{\arraystretch}{1.35}
\begin{longtable}{@{}p{0.33\textwidth}p{0.36\textwidth}p{0.25\textwidth}@{}}
\toprule
\textbf{Formula} & \textbf{Meaning and variables} & \textbf{Conditions / when to use}\\
\midrule
\endhead
\multicolumn{3}{@{}l}{\sffamily\bfseries\color{accent} Logic and sets}\\
$P\Leftrightarrow Q$ iff $(P\Rightarrow Q)$ and $(Q\Rightarrow P)$ & Equivalence is implication both ways & Proving ``iff''; solving\\
$\text{not}(\forall x\in A,P)\Leftrightarrow\exists x\in A:\text{not }P$ & Negating a universal claim & Counterexamples\\
$\text{not}(\exists x\in A:P)\Leftrightarrow\forall x\in A,\text{not }P$ & Negating an existential claim & Contradiction proofs\\
$A=B\Leftrightarrow(A\subseteq B\text{ and }B\subseteq A)$ & Set equality via two inclusions & Proving sets equal\\
$A\setminus B=\{x\in A\mid x\notin B\}$ & Set difference & Definition questions\\
$A\subseteq B\Rightarrow A\setminus B=\emptyset$ & Removing a superset leaves nothing & (Converse also true)\\
\multicolumn{3}{@{}l}{\sffamily\bfseries\color{accent} Real numbers and inequalities}\\
$a>b,\ c>0\Rightarrow ac>bc$;\ \ $a>b,\ c<0\Rightarrow ac<bc$ & M6: sign of multiplier decides direction & Every multiplication/division\\
$a>b\Leftrightarrow a+c>b+c$ & Adding preserves order (A5 twice) & Moving terms across\\
$ax+b>0$: $x>-\frac ba$ ($a>0$), $x<-\frac ba$ ($a<0$) & Linear inequality & Case $a=0$ separately\\
$uv>0\Leftrightarrow$ same signs; $uv<0\Leftrightarrow$ opposite & Sign of a product & Sign analysis\\
$D=b^2-4ac$,\ \ $x_{1,2}=\frac{-b\pm\sqrt D}{2a}$ & Discriminant and roots of $ax^2+bx+c$ & $a\ne0$; real roots iff $D\ge0$\\
$ax^2+bx+c=a(x-x_1)(x-x_2)$ & Factorisation through roots & $D\ge0$\\
$ax^2+bx+c=a\big(x+\frac b{2a}\big)^2-\frac D{4a}$ & Completing the square & $D\le0$; minimum value\\
$a>0$: quadratic $>0$ outside roots, $<0$ between & Sign of a quadratic & Quadratic inequalities\\
\multicolumn{3}{@{}l}{\sffamily\bfseries\color{accent} Absolute value}\\
$|x|=x$ ($x\ge0$), $-x$ ($x<0$);\ $|x|=\max\{x,-x\}$ & Size of $x$ & Definition\\
$-|a|\le a\le|a|$,\ $|{-a}|=|a|$,\ $|ab|=|a||b|$ & Basic properties & Simplifying\\
$|x-z|\le|x-y|+|y-z|$;\ $|a+b|\le|a|+|b|$ & Triangle inequality & Bounding distances; limits\\
$|x-a|<r\Leftrightarrow a-r<x<a+r$ & Within distance $r$ of $a$ & $r>0$\\
\multicolumn{3}{@{}l}{\sffamily\bfseries\color{accent} Sums and induction}\\
$\sum_{j=k}^na_j=a_k+\dots+a_n$ & $n-k+1$ terms & Integers $k\le n$\\
$\sum_{j=k}^{k-1}a_j=0$,\ $\sum_{j=k}^{k}a_j=a_k$ & Empty and one-term sums & Conventions\\
$\sum_{j=k}^na=(n-k+1)a$ & Constant summand & Count terms carefully\\
P1, P2, P3 & Split sums; constants out; join ranges & Manipulating sums\\
$\sum_{j=1}^nj=\frac{n(n+1)}2$;\ $\sum_{j=1}^nj^2=\frac{n(n+1)(2n+1)}6$ & Standard sums & Closed forms\\
$\sum_{t=0}^{n-1}r^t=\frac{1-r^n}{1-r}$ & Geometric sum & $r\ne1$\\
Base $P_1$ + step $P_n\Rightarrow P_{n+1}$ & Induction & Statements for all $n\in\N$\\
\multicolumn{3}{@{}l}{\sffamily\bfseries\color{accent} Binomial formula}\\
$0!=1$, $k!=(k-1)!\,k$;\ \ $C_n^j=\frac{n!}{(n-j)!\,j!}$ & Factorial; binomial coefficient & $0\le j\le n$\\
$(x+y)^n=\sum_{j=0}^nC_n^jx^{n-j}y^j$ & Expansion of a power of a sum & $n\in\{0\}\cup\N$\\
$C_n^j=C_n^{n-j}$,\ \ $C_n^j+C_n^{j+1}=C_{n+1}^{j+1}$ & Symmetry; Pascal's rule & Computing coefficients\\
\multicolumn{3}{@{}l}{\sffamily\bfseries\color{accent} Functions, polynomials, limits}\\
$p(\bar x)=0\Leftrightarrow p(x)=(x-\bar x)q(x)$ & Factor theorem (L2.1), $\deg q=n-1$ & $\deg p=n\ge1$\\
$\forall\varepsilon>0\ \exists N\in\N\ \forall n>N:\ |a_n-x|<\varepsilon$ & $x$ is a limit of $(a_n)$ & Proving convergence\\
$a_n\ge0$, $a_n\to a\Rightarrow a\ge0$ & Limits preserve $\ge$ & Not $>$\\
\bottomrule
\end{longtable}
}

\section*{Glossary}
\addcontentsline{toc}{section}{Glossary}

\glossentry{Absolute value}{$|x|=x$ if $x\ge0$ and $|x|=-x$ if $x<0$; equivalently $|x|=\max\{x,-x\}$.}{How big a number is,
ignoring its sign; $|x-y|$ is how far apart $x$ and $y$ are.}
\glossentry{Axiom}{A proposition accepted without proof.}{A starting rule of the game; in ECN115, the properties O1--M6 of the
real numbers.}
\glossentry{Binomial coefficient}{$C_n^j=\binom nj=\frac{n!}{(n-j)!\,j!}$ for integers $0\le j\le n$.}{The number in front of
$x^{n-j}y^j$ when $(x+y)^n$ is multiplied out.}
\glossentry{Binomial formula}{$(x+y)^n=\sum_{j=0}^nC_n^jx^{n-j}y^j$.}{A one-line recipe for multiplying out any power of a sum.}
\glossentry{Codomain}{The set $R$ in $f\colon D\to R$, in which the values of $f$ lie.}{Where the outputs are allowed to land.}
\glossentry{Coefficient}{In $f(x)=\sum_{k=0}^na_kx^k$, the numbers $a_0,\dots,a_n$.}{The fixed numbers multiplying the powers of $x$.}
\glossentry{Completing the square}{Writing $ax^2+bx+c=a(x+\frac b{2a})^2-\frac D{4a}$.}{Rewriting a quadratic as a square plus a
constant, so its sign and minimum are visible.}
\glossentry{Contrapositive}{Of $P\Rightarrow Q$: the implication $(\text{not }Q)\Rightarrow(\text{not }P)$.}{The same promise read
backwards; always equivalent to the original.}
\glossentry{Converge}{A sequence converges if it has a finite limit.}{It settles down towards a single number.}
\glossentry{Converse}{Of $P\Rightarrow Q$: the implication $Q\Rightarrow P$.}{The reversed claim; it may be false even when the
original is true.}
\glossentry{Counterexample}{An element for which a claimed $\forall$ statement fails.}{One case that breaks a general claim.}
\glossentry{Critical point (of a sign analysis)}{A value at which a factor of the numerator or denominator is zero.}{A place
where an expression can change sign.}
\glossentry{Degree}{For a polynomial with at least one non-zero coefficient, the largest $k$ with $a_k\ne0$.}{The highest
power of $x$ that actually appears.}
\glossentry{Discriminant}{$D=b^2-4ac$ for $ax^2+bx+c$.}{Tells you whether a quadratic has two, one or no real roots.}
\glossentry{Domain}{The set $D$ in $f\colon D\to R$ of inputs.}{The values the function accepts.}
\glossentry{Element}{An object in a set; written $x\in A$.}{One of the things in the bag.}
\glossentry{Empty set}{The set with no elements, $\emptyset$ (or $\{\}$).}{An empty bag.}
\glossentry{Empty sum}{$\sum_{j=k}^{k-1}a_j$, defined to be $0$.}{Adding up nothing gives nothing.}
\glossentry{Equal functions}{Functions with the same domain, the same codomain and the same value at every point.}{Not just the
same formula: the same inputs, outputs and allowed outputs.}
\glossentry{Equal sets}{Sets with exactly the same elements; equivalently, each is a subset of the other.}{Same contents,
however described.}
\glossentry{Equivalence}{$P\Leftrightarrow Q$: $P\Rightarrow Q$ and $Q\Rightarrow P$.}{Two statements that are true in exactly
the same situations.}
\glossentry{Existential quantifier}{$\exists x\in A: P$ --- there is at least one $x\in A$ with property $P$.}{``Some'' or ``there
is''; proved by giving an example.}
\glossentry{Factorial}{$0!=1$ and $k!=(k-1)!\cdot k$.}{$1\times2\times\dots\times k$.}
\glossentry{Function}{A rule (formally, a set of ordered pairs) assigning to each element of a domain a unique element of a
codomain.}{A machine that takes each allowed input to exactly one output.}
\glossentry{Implication}{$P\Rightarrow Q$: if $P$, then $Q$.}{A one-way promise; broken only if $P$ happens and $Q$ does not.}
\glossentry{Induction (mathematical)}{Proving $P_n$ for all $n\in\N$ by proving $P_1$ (base) and $P_n\Rightarrow P_{n+1}$ for all $n$
(step).}{Knock over the first domino and make sure each knocks over the next.}
\glossentry{Induction hypothesis}{The assumption that $P_n$ (or $P_1,\dots,P_n$) holds, made while proving the step.}{What you are
allowed to use about the previous case.}
\glossentry{Integers}{$\Z=\{\dots,-2,-1,0,1,2,\dots\}$.}{Whole numbers, positive, negative and zero.}
\glossentry{Intersection}{$A\cap B$: the objects in both $A$ and $B$.}{What two sets have in common.}
\glossentry{Interval}{A set of real numbers between two endpoints, e.g.\ $[a,b]=\{x\in\R\mid a\le x\le b\}$.}{A stretch of the
number line; square brackets include the endpoint, round ones exclude it.}
\glossentry{Least upper bound}{An upper bound $z$ of $A$ such that every $w<z$ is exceeded by some element of $A$.}{The
smallest number that is at least every element of the set (non-examinable).}
\glossentry{Limit (of a sequence)}{$x$ such that $\forall\varepsilon>0\ \exists N\in\N\ \forall n>N:\ |a_n-x|<\varepsilon$.}{The
number the terms eventually stay as close to as you like.}
\glossentry{Monomial}{A single term $a_kx^k$ of a polynomial.}{One piece of a polynomial.}
\glossentry{Natural numbers}{$\N=\{1,2,3,\dots\}$ (not including $0$ in this module).}{The counting numbers.}
\glossentry{Negation}{``not $P$'': true exactly when $P$ is false.}{The opposite claim.}
\glossentry{Numerical function}{A function $f\colon D\to R$ with $D,R\subseteq\R$.}{A function from numbers to numbers.}
\glossentry{Numerical sequence}{A sequence whose terms are real numbers.}{An endless list of numbers.}
\glossentry{Parameter}{A letter standing for a fixed but unspecified number.}{A number in the problem that you are not told,
so the answer depends on it.}
\glossentry{Polynomial}{A function $f(x)=\sum_{k=0}^na_kx^k$.}{A sum of multiples of powers of $x$.}
\glossentry{Proof}{A demonstration that one proposition follows from one or more others.}{A step-by-step argument that leaves no
room for doubt.}
\glossentry{Proposition}{A statement that may be either true or false.}{A claim.}
\glossentry{Rational numbers}{$\Q$: numbers $p/q$ with $p\in\Z$, $q\in\N$.}{Fractions.}
\glossentry{Real numbers}{$\R$: informally, the points of the number line.}{All the numbers on the line, including $\sqrt2$ and $\pi$.}
\glossentry{Root}{A value $\bar x$ with $f(\bar x)=0$.}{Where the graph meets the horizontal axis.}
\glossentry{Sequence}{An infinite ordered list $(x_1,x_2,\dots)$; formally a function with domain $\N$.}{A list of values,
one for each period $1,2,3,\dots$.}
\glossentry{Set}{A collection of objects (taken as a basic notion).}{A bag of things; only membership matters.}
\glossentry{Set-builder notation}{$\{x\in A\mid x\text{ has property }P\}$.}{``The members of $A$ that satisfy $P$.''}
\glossentry{Set difference}{$A\setminus B=\{x\in A\mid x\notin B\}$.}{$A$ with everything in $B$ thrown out.}
\glossentry{Solution set}{The set of all values for which an equation or inequality holds.}{Every answer, collected together.}
\glossentry{Subset}{$A\subseteq B$: every element of $A$ is an element of $B$.}{$A$ fits inside $B$.}
\glossentry{Successor property}{For any $n\in\N$, $n+1\in\N$.}{There is always a next natural number.}
\glossentry{Summation operator}{$\sum_{j=k}^na_j=a_k+a_{k+1}+\dots+a_n$.}{Shorthand for adding a list of terms.}
\glossentry{Superset}{$B$ is a superset of $A$ if $A\subseteq B$.}{The bigger set that contains $A$.}
\glossentry{Theorem}{A proposition proved on the basis of axioms and previously established theorems.}{A claim that has been
shown to follow from the rules.}
\glossentry{Triangle inequality}{$|x-z|\le|x-y|+|y-z|$; equivalently $|a+b|\le|a|+|b|$.}{A detour is never shorter than the
direct route.}
\glossentry{Trichotomy}{Property O1: for real $a,b$ exactly one of $a>b$, $a=b$, $a<b$ holds.}{Any two numbers can be compared,
in exactly one way.}
\glossentry{Union}{$A\cup B$: the objects in $A$, in $B$, or in both.}{Everything in either set.}
\glossentry{Universal quantifier}{$\forall x\in A,\ P$ --- every $x\in A$ has property $P$.}{``All''; proved for an arbitrary
element.}
\glossentry{Upper bound}{$y$ with $x\le y$ for every $x\in A$.}{A ceiling for the set.}

\section*{What do I actually need to know?}
\addcontentsline{toc}{section}{What do I actually need to know?}

Tick each item only when you can do it without notes.

\subsection*{Chapter 1: The language of mathematics}
\begin{checklist}
  \item Define proposition, proof, axiom and theorem.
  \item Explain $P\Rightarrow Q$, $P\Leftrightarrow Q$ and ``not $P$''; give the converse and contrapositive of an implication.
  \item Negate inequalities, ``and''/``or'' statements and quantified statements.
  \item Prove $\exists$ statements by example and disprove $\forall$ statements by counterexample.
  \item Explain why $\forall n\,\exists m$ and $\exists m\,\forall n$ differ.
\end{checklist}
\subsection*{Chapter 2: Sets}
\begin{checklist}
  \item Define subset, set equality, intersection, union, set difference and set-builder notation.
  \item Calculate $A\cap B$, $A\cup B$, $A\setminus B$ for given sets.
  \item Prove a subset relation for arbitrary sets.
  \item Answer ``find all sets $D$ such that\dots'' and ``are there sets such that\dots'' questions.
\end{checklist}
\subsection*{Chapter 3: Number systems}
\begin{checklist}
  \item State $\N,\Z,\Q,\R$, their inclusions and the successor property.
  \item State O1--O2, A1--A5, M1--M6, especially both clauses of M6.
  \item Explain why $\R$ is needed beyond $\Q$.
  \item Write intervals in set-builder notation and back.
\end{checklist}
\subsection*{Chapter 4: Inequalities}
\begin{checklist}
  \item Solve linear inequalities, reversing the direction when dividing by a negative number.
  \item Solve parameter inequalities by cases $a>0$, $a=0$, $a<0$.
  \item Solve product, quadratic and fractional inequalities by sign analysis.
  \item Use the discriminant and completing the square.
\end{checklist}
\subsection*{Chapter 5: Absolute value}
\begin{checklist}
  \item Define $|x|$ both ways and prove its basic properties.
  \item State and apply the triangle inequality; derive $|a+b|\le|a|+|b|$ and $|a|\le|a-b|+|b|$.
  \item Solve $|x-a|<r$ and $|x-a|>r$.
\end{checklist}
\subsection*{Chapter 6: Summation}
\begin{checklist}
  \item Evaluate sums with any integer limits; count $n-k+1$ terms.
  \item State the one-term and empty-sum conventions and P1--P3.
  \item Find closed forms using P1, P2 and the constant-sum rule.
\end{checklist}
\subsection*{Chapter 7: Induction}
\begin{checklist}
  \item State the principle of induction with base, step and hypothesis.
  \item Prove summation formulas and simple inequalities by induction.
  \item State the alternative (strong) formulation and use it with several base cases.
\end{checklist}
\subsection*{Chapter 8: Binomial formula}
\begin{checklist}
  \item Compute $k!$ and $C_n^j$; state the binomial formula.
  \item Expand powers and find specific coefficients.
  \item Prove lower bounds such as $1.02^{102}>6$ using partial sums.
\end{checklist}
\subsection*{Chapter 9: Functions and polynomials}
\begin{checklist}
  \item Give the informal and formal definitions of a function and of equal functions.
  \item Define polynomial, degree, monomial and root.
  \item State Theorem L2.1 and use it to factorise polynomials.
\end{checklist}
\subsection*{Chapter 10: Sequences and limits}
\begin{checklist}
  \item Define sequence and numerical sequence; state the definition of a limit exactly.
  \item Find $N$ for a given $\varepsilon$; prove convergence from the definition.
  \item Prove that a sequence such as $1,0,1,0,\dots$ has no limit.
  \item Prove that limits preserve non-strict inequalities.
\end{checklist}

\section*{Final practice paper}
\addcontentsline{toc}{section}{Final practice paper}

This paper is written in the style of the past ECN115 papers (definitions, ``if yes give an example, if no argue why'',
parameter questions and ``explain why the argument is wrong''), but uses only the material in this book. Total: 100 marks.

\setcounter{pq}{0}
\renewcommand{\thepq}{F.\arabic{pq}}

\pq[18 marks] Let $A$, $B$ be sets.
(a) Define the set difference $A\setminus B$ and state what it means for $A$ and $B$ to be equal in terms of $\subseteq$. [5]
(b) Are there non-empty sets of natural numbers with $A\cup B=A$ and $A\ne B$? If yes, give an example; if no, argue why. [4]
(c) Find all sets $D\subseteq\N$ such that $\{1,2\}\cup D=\{1,2,3,4\}$ and $1\notin D$. [5]
(d) Prove that $A\setminus B\subseteq A$ for all sets $A,B$. [4]

\pq[22 marks]
(a) Solve $4-\tfrac{3}{2}x\ge x+9$, naming the property that decides the direction at the step where you divide. [5]
(b) Solve $(2-x)(3x+1)<0$. [5]
(c) Let $a\in\R$ be a parameter. Find all $x$ with $(a-2)x\le a-2$. Your answer should depend on $a$. [6]
(d) Solve $\dfrac{x+2}{x-1}\le 2$. [6]

\pq[16 marks]
(a) State the triangle inequality for real numbers $x,y,z$. [2]
(b) Prove that $|a|-|b|\le|a-b|$ for all real $a,b$. [4]
(c) Solve $|3x-2|\le 4$ and $|x+1|\ge|x-3|$. [6]
(d) Show that if $|x-2|<1$ then $|x^2-4|<5$. [4]

\pq[16 marks]
(a) Evaluate $\sum_{j=-2}^{3}(2j+1)$ and $\sum_{j=4}^{3}j^5$. [3]
(b) Prove by induction that $\sum_{j=1}^{n}(3j-1)=\frac{n(3n+1)}{2}$ for all $n\in\N$, identifying the base, the hypothesis and
where it is used. [8]
(c) A student checks the identity in (b) for $n=1,\dots,20$ and says this proves it. Explain why this is not a proof. [2]
(d) Find $\sum_{j=1}^{n}(3j-1)$ using P1, P2 and $\sum_{j=1}^nj=\frac{n(n+1)}2$ instead. [3]

\pq[14 marks]
(a) Write the binomial formula for $(1+x)^n$ and the first four terms explicitly. [5]
(b) Show that $1.05^{40}>6$. [6]
(c) Find the coefficient of $x^4$ in $(1-2x)^7$. [3]

\pq[14 marks]
(a) Define what it means that a number $x$ is a limit of a sequence $(a_n)$. [4]
(b) Prove from the definition that $\frac{2n+3}{n+1}\to2$. [5]
(c) A student argues: ``The sequence $a_n=(-1)^n\frac{n}{n+1}$ has terms that get closer and closer to $1$ and to $-1$, so it
converges to both $1$ and $-1$.'' Explain why this argument is wrong and prove that $(a_n)$ has no limit. [5]

\section*{Solutions to the final practice paper}
\addcontentsline{toc}{section}{Solutions to the final practice paper}

\pqsol{F.1}
(a) $A\setminus B=\{x\in A\mid x\notin B\}$, the elements of $A$ that are not elements of $B$. $A=B$ if and only if
$A\subseteq B$ and $B\subseteq A$.
(b) Yes. $A\cup B=A$ holds exactly when $B\subseteq A$. Take $A=\{1,2\}$, $B=\{1\}$: both non-empty, $A\cup B=\{1,2\}=A$ and
$A\ne B$.
(c) $3,4\in D$ (they are in the union but not in $\{1,2\}$); $D\subseteq\{1,2,3,4\}$; $1\notin D$; $2$ is free. So
$D\in\{\{3,4\},\{2,3,4\}\}$.
(d) Let $x\in A\setminus B$. By definition $x\in A$ and $x\notin B$; in particular $x\in A$. As $x$ was arbitrary,
$A\setminus B\subseteq A$.

\pqsol{F.2}
(a) $4-\tfrac32x\ge x+9\Leftrightarrow-\tfrac52x\ge5$. Multiplying by $-\tfrac25<0$ reverses the inequality (M6, second
clause): $x\le-2$. Solution $(-\infty,-2]$.
(b) Critical points $x=2$ and $x=-\tfrac13$. The product is negative when the signs differ: $2-x>0$ and $3x+1<0$ gives
$x<-\tfrac13$; $2-x<0$ and $3x+1>0$ gives $x>2$. Solution $(-\infty,-\tfrac13)\cup(2,\infty)$.
(c) If $a>2$: divide by $a-2>0$: $x\le1$, i.e.\ $(-\infty,1]$. If $a<2$: divide by $a-2<0$, reversing: $x\ge1$, i.e.\
$[1,\infty)$. If $a=2$: $0\le0$, true for all $x$: $\R$.
(d) $\frac{x+2}{x-1}-2\le0\Leftrightarrow\frac{x+2-2x+2}{x-1}\le0\Leftrightarrow\frac{4-x}{x-1}\le0$. Critical points $1$ (excluded)
and $4$ (included). For $x<1$: numerator $+$, denominator $-$: quotient $-$, included. For $1<x<4$: both $+$: excluded. For
$x>4$: numerator $-$, denominator $+$: quotient $-$, included. Solution $(-\infty,1)\cup[4,\infty)$.

\pqsol{F.3}
(a) $|x-z|\le|x-y|+|y-z|$.
(b) With $x=a$, $y=b$, $z=0$: $|a|\le|a-b|+|b|$, so $|a|-|b|\le|a-b|$.
(c) $|3x-2|\le4\Leftrightarrow-4\le3x-2\le4\Leftrightarrow-\tfrac23\le x\le2$: $[-\tfrac23,2]$. For the second, square (both sides
$\ge0$): $(x+1)^2\ge(x-3)^2\Leftrightarrow2x+1\ge-6x+9\Leftrightarrow8x\ge8\Leftrightarrow x\ge1$: $[1,\infty)$ (points at least as far
from $-1$ as from $3$).
(d) If $|x-2|<1$ then $1<x<3$, so $3<x+2<5$ and $|x+2|<5$; hence $|x^2-4|=|x-2|\,|x+2|<1\cdot5=5$.

\pqsol{F.4}
(a) $-3-1+1+3+5+7=12$; the second is an empty sum, $0$.
(b) Let $P_n$ be the identity. Base: $3-1=2=\frac{1\cdot4}{2}$. Hypothesis: assume $P_n$. Step:
$\sum_{j=1}^{n+1}(3j-1)=\frac{n(3n+1)}2+(3n+2)$ (hypothesis used) $=\frac{3n^2+7n+4}{2}=\frac{(n+1)(3n+4)}{2}=\frac{(n+1)(3(n+1)+1)}{2}$,
which is $P_{n+1}$. By induction $P_n$ holds for all $n\in\N$.
(c) It verifies twenty propositions $P_1,\dots,P_{20}$; the claim concerns infinitely many, and nothing is said about
$P_{21}$ onwards. A proof needs the general step.
(d) $3\sum j-\sum1=\frac{3n(n+1)}2-n=\frac{3n^2+n}{2}=\frac{n(3n+1)}2$.

\pqsol{F.5}
(a) $(1+x)^n=\sum_{j=0}^nC_n^jx^j=1+nx+\frac{n(n-1)}2x^2+\frac{n(n-1)(n-2)}6x^3+\dots$
(b) With $x=0.05$, $n=40$, all terms positive:
$1.05^{40}>1+40(0.05)+780(0.0025)+9880(0.000125)=1+2+1.95+1.235=6.185>6$
(using $C_{40}^2=780$, $C_{40}^3=9880$).
(c) General term $C_7^j(-2x)^j$; for $j=4$: $C_7^4\cdot16=35\cdot16=560$.

\pqsol{F.6}
(a) $x$ is a limit of $(a_n)$ if for every $\varepsilon>0$ there exists $N\in\N$ such that $|a_n-x|<\varepsilon$ for all $n>N$.
(b) $\big|\frac{2n+3}{n+1}-2\big|=\frac{1}{n+1}<\frac1n$. Given $\varepsilon>0$, choose $N\ge\frac1\varepsilon$. For $n>N$:
$\big|\frac{2n+3}{n+1}-2\big|<\frac1n<\frac1N\le\varepsilon$.
(c) A sequence has at most one limit, and being close to two values infinitely often prevents convergence rather than
giving two limits. Proof: suppose $x$ is a limit, and take $\varepsilon=\tfrac12$. There is $N$ with $|a_n-x|<\tfrac12$ for $n>N$.
Choose an even $m>N$ and an odd $k>N$ with $m,k\ge3$; then $a_m=\frac m{m+1}\ge\tfrac34$ and $a_k=-\frac{k}{k+1}\le-\tfrac34$. By the
triangle inequality, $\tfrac32\le|a_m-a_k|\le|a_m-x|+|x-a_k|<\tfrac12+\tfrac12=1$, a contradiction. So $(a_n)$ has no limit.
